% Laget automatisk av elevark/lag_elevark.py fra sim/lysavboyning.html. Endringer her blir overskrevet.
\documentclass[11pt]{article}
\usepackage{elevark}
\begin{document}
\elevarkhode{Fysikk 2 · Generell relativitetsteori}{Lys som bøyes av tyngdekraften}{https://fysikklab.no/sim/lysavboyning.html}{Simuleringen}

Ifølge generell relativitetsteori krummer masse rommet, og lys som går forbi en stjerne, får banen sin bøyd. Under solformørkelsen i 1919 målte Arthur Eddington at stjerner nær solranden så ut til å være flyttet litt, akkurat så mye som Einstein hadde forutsagt. Hvor stor blir avbøyningen ved andre himmellegemer?

\section*{Oppgaver}

\begin{oppgaveliste}
\oppgave \fase{gjett}{Gjett}En stjerne står nesten bak sola. Ser den ut til å være nærmere sola eller lenger unna enn den egentlig er?
\svar
\oppgave \fase{test}{Regn}Regn ut avbøyningen for lys som streifer solranden ($b = 6{,}96 \cdot 10^{8}$ $m$). Sjekk med simuleringen. Hvor stor blir den i dobbel avstand?
\svar
\oppgave \fase{forklar}{Forklar}Hvorfor kunne Eddington bare måle dette under en total solformørkelse? Hvorfor var det viktig at resultatet var 1,75″ og ikke 0,87″?
\svar
\oppgave \fase{test}{Undersøk}Sammenlign avbøyningen ved overflaten av jorda, sola, en hvit dverg og en nøytronstjerne. Hva er det som avgjør hvor stor avbøyningen blir?
\svar
\oppgave \fase{test}{Regn}Regn ut Schwarzschild-radien til sola og til jorda. Hva ville skjedd om sola ble presset sammen til en kule med denne radien?
\svar
\oppgave \fase{utfordring}{Utfordring}Lys fra en fjern galakse går forbi en galaksehop på begge sider og kommer fram til oss. Hva ser vi? Søk opp «gravitasjonslinse» eller «Einsteinring».
\svar
\end{oppgaveliste}

\end{document}
